\documentclass[a4paper]{article}
\usepackage{amsfonts}
\usepackage{amsmath}
\usepackage{amssymb}
\usepackage{graphicx}     

\begin{document}
  
\noindent \large Auteur: Anna Voronovska \\
\noindent \large Studierichting: Informatica\\
\noindent \large Oefeningenreeks 7A nr. 11.2\\

\medskip

\normalsize

$\mathbf{f} : \Omega \subseteq \mathbb{R} \longrightarrow \mathbb{R}^4 : x \mapsto (2^{\sec x},\ \coth x,\ \log_5\sqrt{\cos x},\ x^x)$\\

$D\mathbf{f}(x) = \left(\begin{array}{c} \dfrac{\partial f_1}{\partial x} \\ \dfrac{\partial f_2}{\partial x} \\ \dfrac{\partial f_3}{\partial x} \\ \dfrac{\partial f_4}{\partial x} \end{array}\right)$\\

\begin{itemize}

\item $\dfrac{\partial f_1}{\partial x} = 2^{\sec x} \ln 2 \cdot \sec x \tan x$\\

\item $\dfrac{\partial f_2}{\partial x} = -\dfrac{1}{\sinh^2 x}$\\

\item $f_3 = \log_5\sqrt{\cos x} = \dfrac{\ln(\cos x)}{2\ln 5}$\\

$\dfrac{\partial f_3}{\partial x} = \dfrac{1}{2\ln 5} \cdot \dfrac{-\sin x}{\cos x} = -\dfrac{\tan x}{2\ln 5}$\\

\item $f_4 = x^x = e^{x\ln x}$\\

$\dfrac{\partial f_4}{\partial x} = x^x(\ln x + 1)$\\

\end{itemize}

$\Rightarrow D\mathbf{f}(x) = \left(\begin{array}{c} 2^{\sec x} \ln 2 \cdot \sec x \tan x \\ -\dfrac{1}{\sinh^2 x} \\
-\dfrac{\tan x}{2\ln 5} \\ x^x(\ln x + 1) \end{array}\right)$\\

\begin{thebibliography}{999}
%\bibitem{a} Referentie
\end{thebibliography}
\end{document} 
